% !TEX root = ../main.tex
\PlanPage{自研成果与持续积累}{武桐宇}{35}
\textbf{本页结论：}项目当前可展示的积累是可运行的样例流程、结构化事件/图谱资产与确定性评分实现；外部政策、公告和未来授权财务数据属于引用或待接入资源，不能写为自研数据资产。

自研边界应与复用边界同时交代。事件字段、图谱文件、规则代码、导出约定和前端展示共同组成原型；政策公告、监管信息和公司材料则必须保留原始链接和来源标识。这样既能说明团队投入，也不把公共信息重新包装成独占内容。

\begin{table}[htbp]
\centering
\caption{项目资产的自研、引用与待接入边界}
\label{tab:wty-own-work}
\small
\begin{tabularx}{\linewidth}{p{0.23\linewidth}X p{0.24\linewidth}}
\toprule
资产类别 & 当前工作内容 & 状态与边界 \\
\midrule
事件结构与证据规范 & 定义事件字段、来源状态与核验流程 & 自研结构；原始事实仍来自公开来源 \\
产业图谱包 & 节点、关系、公司映射和事件规则的文件化管理 & 自研组织；部分关系/规则待核验 \\
评分实现 & $E\times S\times F/10000$ 的确定性计算与方向字段 & 自研实现；不等于金融预测模型 \\
界面与导出流程 & 对样例事件、路径和分项进行展示与导出 & 当前默认 \texttt{mock-v1} 演示 \\
授权财务数据接入 & 字段字典与清洗接口预留 & 尚未作为正式数据接入或发布 \\
\bottomrule
\end{tabularx}
\end{table}

持续积累的重点是为每次新增关系、新来源和新规则留下版本差异，而非无限扩张概念标签。任何外部数据、开源组件或受限资料在使用前都应核对许可与再分发条件。
